% TOP_PROBLEM: {"id": 3130, "title": "Signing a graph to have small magnitude eigenvalues", "category_id": 3, "category": {"id": 3, "name": "graph_theory", "display_name": "Graph Theory", "description": "Problems involving graphs, networks, and their properties.", "slug": "graph-theory", "order_index": 3, "created_at": "2026-07-31T15:26:25.671Z"}, "status": "open", "problem_number": "OPG-48264", "set_id": 12, "statusQualification": "The necessary degree range d at least two is explicit; d=1 would give an immediate counterexample to the unqualified shorthand."}
% FIRST_POSTED: 2026-10-01
% ENTRY_KIND: statement_only
% UnsolvedMath ID 3130; source catalogue code OPG-48264
% !TeX program = lualatex
% Definition and sources only; this file is not an LLM solution attempt.
\documentclass[11pt,a4paper]{article}
\usepackage[margin=25mm]{geometry}
\usepackage{fontspec}
\setmainfont{lmroman10-regular.otf}[BoldFont=lmroman10-bold.otf,
 ItalicFont=lmroman10-italic.otf,BoldItalicFont=lmroman10-bolditalic.otf]
\usepackage{amsmath,amssymb,mathtools}
\usepackage{unicode-math}
\setmathfont{latinmodern-math.otf}
\AtBeginDocument{\let\setminus\smallsetminus}
\usepackage{xurl,hyperref}
\hypersetup{colorlinks=true,urlcolor=blue}
\setcounter{secnumdepth}{0}
\setlength{\parindent}{0pt}
\setlength{\parskip}{5pt}
\setlength{\emergencystretch}{3em}
\begin{document}
\section*{Signing a graph to have small magnitude eigenvalues}\label{3130}
{\small UnsolvedMath ID 3130 · OPG-48264 · Graph Theory · OpenGarden}
\subsection{Definitions and mathematical statement}
Let $d\ge2$ and let $G=(V,E)$ be a finite simple undirected $d$-regular
graph, meaning that every vertex has exactly $d$ neighbours.
For a signing $s:E\to\{-1,1\}$ define the real symmetric matrix
$A_s$, indexed by $V$, by
\[
 (A_s)_{uv}=\begin{cases}s(\{u,v\}),&\{u,v\}\in E,\\
                         0,&\text{otherwise}.\end{cases}
\]
In particular its diagonal is zero, and the two entries corresponding
to the same undirected edge receive the same sign. A signing does not
change which entries of the adjacency matrix are zero.

The conjecture asserts that for every such $G$ there is a signing with
\[
 \max_{\lambda\in\operatorname{Spec}(A_s)}|\lambda|
                 \le2\sqrt{d-1}.
\]
All eigenvalues are real because $A_s$ is symmetric. The inequality
bounds both the largest eigenvalue and the magnitude of the most
negative one; a bound on only one end of the spectrum is insufficient.
Equivalently, the Euclidean operator norm of $A_s$ must have this bound.

Every eigenvalue of the signed matrix is included, with no exception
for eigenvectors related to those of the unsigned regular graph.
The signs may depend on the whole graph and need not come from
assigning signs to vertices. No bipartiteness or connectedness is imposed.
The source's implicit range is specified as $d\ge2$: at $d=1$ a single
edge has signed eigenvalues $1,-1$, whereas the proposed right-hand
side is zero, so that degenerate case cannot be included.
\subsection{Short English statement}
Can every regular graph of degree at least two be signed so that all signed adjacency eigenvalues have magnitude at most $2\sqrt{d-1}$?
\subsection{Sources}
\begin{itemize}
\item[S1] ULAM.AI and the UnsolvedMath Contributors, supplied problems.json, record 3130 (OPG-48264), OpenGarden collection. Catalogue identity and source transcription; CC BY 4.0.
\url{https://huggingface.co/datasets/ulamai/UnsolvedMath}.
\item[S2] Open Problem Garden, Signing a graph to have small magnitude eigenvalues. Original problem and discussion; the mathematical formulation below is expanded from the supplied source record.
\url{http://www.openproblemgarden.org/op/signing_a_graph_to_have_small_magnitude_eigenvalues}.
\item[S3] Y. Bilu and N. Linial, Constructing expander graphs by 2-lifts and discrepancy vs. spectral gap, original signing framework, arXiv:math/0312022.
\url{https://arxiv.org/abs/math/0312022}.
\end{itemize}
\end{document}
